\begin{textsample}{NEG Dim 1 – 2014\_c6 – Score -9.00 – 2014\_c6/enade\_2014\_tec\_autom\_43f2ddc563.txt}  \label{ex:f1_neg_008}
CAPÍTULO 3 ANÁLISE TÉCNICA DA PROVA Este capítulo tem por objetivo apresentar o desempenho dos estudantes concluintes de Tecnologia em Automação Industrial no ENADE/2014.
Para isso, foram calculadas as estatísticas básicas da prova em seu todo ( seção 3.1.1 ), bem como as estatísticas relacionadas aos Componentes de Formação Geral ( seção 3.1.2 ) e de Conhecimento Específico da Área ( seção 3.1.3 ).
Dadas as suas características, foram analisadas, em separado, as questões objetivas ( seção 3.2 ) e as questões discursivas ( seção 3.3 ).
Para as questões objetivas, foram disponibilizados os índices de facilidade e de discriminação bisserial, também em separado, para os Componentes de Formação Geral ( seção 3.2.1 ) e de Conhecimento Específico ( seção 3.2.2 ).
De \textbf{cada} componente, uma das questões foi escolhida para exemplificar a análise gráfica, relacionando as alternativas escolhidas pelos estudantes ( inclusive o gabarito ) com o número de acertos no componente.
O Anexo I apresenta a íntegra da análise gráfica para todas as questões objetivas.
Para \textbf{cada} uma das questões discursivas, os conteúdos dos tipos mais comuns de respostas dos estudantes são apresentados e comparados com o padrão de respostas esperado ( ver Anexo VIII com o padrão de respostas ).
Tomando como base as duas questões discursivas do Componente de Formação Geral, a seção 3.3.1.6 apresenta comentários sobre a correção das respostas com respeito à Língua Portuguesa.
Na Tabela 3.1 são apresentados o tamanho da população inscrita e de presentes, e as seguintes estatísticas das notas18 : média do desempenho na prova, erro padrão da média, desvio padrão, nota mínima, mediana e nota máxima.
Tais estatísticas contemplam o total de estudantes concluintes da Área de Tecnologia em Automação Industrial inscritos e presentes à prova do ENADE/2014, tendo em vista as agregações por Grandes Regiões e o país como um todo.
As estatísticas das demais tabelas ímpares serão baseadas na mesma população da Tabela 3.1 e, portanto, do mesmo tamanho.
As três primeiras tabelas pares são desagregadas por Categoria Administrativa e Organização Acadêmica da IES.
As estatísticas dessas tabelas pares serão baseadas na mesma população da Tabela 3.2.
Em relação aos gráficos de distribuição de notas, o intervalo considerado foi de 10 unidades, aberto à esquerda e fechado à direita, com exceção do primeiro intervalo, [ 0 ; 10 ], fechado em ambos os extremos.
Para os gráficos de distribuição das notas das questões discursivas, foram consideradas mais duas categorias : questão em branco19 e nota zero. 18 Uma definição dessas estatísticas pode ser encontrada no Glossário. notas das questões discursivas, foram consideradas mais duas categorias : questão em branco19 e nota zero. 3.1 ESTATÍSTICAS BÁSICAS DA PROVA Esta seção apresenta estatísticas selecionadas e histogramas da nota geral ( 3.1.1 ) e de \textbf{cada} componente : Formação Geral ( 3.1.2 ) e Conhecimento Específico ( 3.1.3 ).
São também apresentadas estatísticas selecionadas de subpopulações caracterizadas por Grande Região, Categoria Administrativa e Organização Acadêmica. 3.1.1 Estatísticas Básicas Gerais A Tabela 3.1 apresenta as estatísticas básicas da prova, por grande Região, dos estudantes concluintes de Tecnologia em Automação Industrial.
A população total de inscritos foi de 2.229.
Destes, 1.689 \textbf{estiveram} presentes, sendo 24,2 \% o índice de não comparecimento.
A Região de maior abstenção foi a Sudeste ( 25,8 \%, ) e a de menor abstenção foi a Norte ( 15,4 \% ).
A média das notas da prova como um todo ( nas seções seguintes serão analisados os Componentes de Formação Geral e de Conhecimento Específico ) foi 48,0, sendo que os alunos da região Centro-Oeste obtiveram a média mais baixa ( 44,3 ), e os da região Nordeste obtiveram a média mais alta ( 52,3 ).
As demais médias foram : 47,6 na região Norte ; 47,1 na região Sudeste ; e 50,6 na região Sul.
O desvio padrão para o Brasil como um todo foi 14,5, sendo o maior desvio padrão encontrado na região Norte ( 16,0 ) e o menor, na região Nordeste ( 12,6 ), indicando uma dispersão um pouco menor das notas desta última região.
A região que obteve a maior nota máxima foi a Sudeste ( 84,3 ), ao passo que a região que atingiu a menor nota máxima foi a Norte ( 71,2 ).
A mediana do Brasil como um todo foi 48,2, sendo a maior mediana obtida na região Sul ( 52,2 ), e a menor obtida na Centro-Oeste ( 44,4 ).
A nota mínima foi obtida na região Sudeste ( 6,8 ), enquanto a maior nota mínima foi obtida na região Nordeste ( 25,6 ).
Considerando -se as notas segundo Grande Região, observa -se que existe diferença estatisticamente significativa ao nível de 95 \% entre a menor média, obtida na região Centro-Oeste ( 44,3 ), e a maior média, obtida na região Nordeste ( 52,3 ).
O comportamento das notas dos estudantes de todo o Brasil pode ser observado no Gráfico 3.1 que apresenta um histograma com a distribuição das mesmas.
Essa é uma distribuição unimodal com a moda no intervalo ( 40 ; 50 ].
Destaca -se também, o intervalo ( 50 ; 60 ], com frequência muito próxima da verificada para o intervalo modal.
A Tabela 3.2 apresenta as estatísticas básicas da prova, desagregadas por Categoria Administrativa e por Organização Acadêmica.
Da população total de inscritos, 1.539 são de IES Privadas e 690 de IES Públicas.
Em relação à Organização Acadêmica, a maior participação foi obtida por estudantes de Universidades ( 1.087 ), seguido de Faculdades ( 850 ) e Centros Universitários ( 292 ).
A Categoria Administrativa de maior abstenção foi a Privada ( 25,3 \% ), e entre as Organizações Acadêmicas foi a das Universidades ( 30,4 \% ), ambos acima da média nacional de 24,2 \%.
A média das notas da prova como um todo foi 48,0.
Em relação à Categoria Administrativa, os estudantes das IES Públicas obtiveram média mais alta ( 55,0 ), e os das IES Privadas obtiveram média mais baixa ( 44,6 ), que a média nacional.
Observa -se que existe diferença estatisticamente significativa entre as médias das notas das IES Públicas e Privadas.
A diferença entre as médias das regiões Nordeste e Centro-Oeste ( 8,0 ), a maior e a menor média, é inferior à diferença entre IES Públicas e Privadas ( 10,4 ), caracterizando uma maior diversidade administrativa do que regional.
No tocante à Organização Acadêmica, apenas as Universidades obtiveram média mais alta ( 49,2 ) que a nacional.
As demais médias foram : 43,5 nos Centros Universitários e 48,0 nas Faculdades.
O desvio padrão para as Faculdades ( 14,8 ) foi superior ao do Brasil como um todo ( 14,5 ), indicando uma dispersão um pouco maior das notas nesta Organização Acadêmica.
Constata -se que existe diferença estatisticamente significativa ao nível de 95 \% nas médias das notas dos estudantes provenientes de Centros Universitários em relação aos de Universidades e Faculdades. 3.1.2 Estatísticas Básicas no Componente de Formação Geral A Tabela 3.3 apresenta as estatísticas básicas em relação ao componente da prova que avalia a Formação Geral dos estudantes concluintes.
Os alunos de todo Brasil obtiveram desempenho médio de 53,8.
Quanto à variabilidade, o desvio padrão das notas dos estudantes do Brasil, como um todo, foi 17,4.
A maior média foi obtida na região Nordeste ( 59,8 ), e a menor, na região Centro-Oeste ( 46,9 ).
As demais médias foram : 58,7 na região Norte ; 54,1 na região Sudeste ; e 54,0 na região Sul.
Já o maior desvio padrão foi obtido na região Centro-Oeste ( 19,9 ), e o menor, na região Nordeste ( 14,6 ).
Os demais desvios padrões foram : 17,7 na região Norte ; 17,2 na região Sudeste ; e 17,0 na região Sul.
A maior nota no Componente de Formação Geral da prova do ENADE/2014 foi 98,0, obtida por, pelo menos, um aluno na região Sul.
A menor nota máxima foi obtida na região Norte ( 87,5 ).
A mediana do Brasil, como um todo, foi 55,4, a mesma obtida na região Sul, sendo a menor mediana encontrada na região Centro-Oeste ( 47,0 ) ; e a maior encontrada na região Norte ( 59,9 ).
A nota mínima nesta parte foi zero em quase todas as regiões, exceto nas regiões Norte ( 22,5 ) e Nordeste ( 22,1 ).
Considerando -se as notas segundo Grande Região, observa -se que existe diferença estatisticamente significativa entre a maior média das notas do Componente de Formação Geral, obtida na região Nordeste ( 59,8 ), e na região Centro-Oeste ( 46,9 ), a menor média.
O Gráfico 3.2 propicia a avaliação do desempenho dos estudantes do Componente de Formação Geral, a partir do histograma da distribuição das notas correspondentes.
A distribuição é unimodal, com moda em ( 50 ; 60 ], enquanto na prova, como um todo ( Gráfico 3.1 ), a moda foi alcançada um intervalo abaixo ( 40 ; 50 ].
Nota -se, ainda que, no Gráfico 3.2, as notas apresentam maior dispersão do que no Gráfico 3.1 ( distribuição das notas da prova ), o que se confirma pela comparação dos desvios padrões : 14,5 para a nota da prova como um todo e 17,4 para o Componente de Formação Geral.
Na Tabela 3.4 são apresentadas as informações referentes ao desempenho dos concluintes do Componente de Formação Geral, em dois tipos de agregações : Categoria Administrativa e Organização Acadêmica.
Observa -se que existe diferença estatisticamente significativa entre as médias dos tipos de Categoria Administrativa.
A maior média foi obtida por estudantes de IES Públicas ( 59,8 ), com uma diferença estatisticamente significativa da obtida por estudantes de IES Privadas ( 51,1 ).
Considerando -se o tipo de Organização Acadêmica, nota -se que não há diferença estatisticamente significativa entre as médias dos alunos de Universidades ( 55,8 ) e de Faculdades ( 53,5 ) ; mas a média dos estudantes de Centros Universitários ( 48,2 ) é estatisticamente diferente em relação às médias das outras Organizações Acadêmicas. 3.1.3 Estatísticas Básicas do Componente de Conhecimento Específico A Tabela 3.5 apresenta as estatísticas básicas referentes ao Componente de Conhecimento Específico da Área de Tecnologia em Automação Industrial.
A média do desempenho dos alunos do Brasil, como um todo, foi 46,0.
A maior média foi obtida na região Nordeste ( 49,7 ), e a menor, na região Centro-Oeste ( 43,5 ).
As demais médias foram : 43,9 na região Norte ; 44,8 na região Sudeste ; e 49,5 na região Sul.
Quanto à variabilidade das notas, o desvio padrão do Brasil, como um todo, foi 16,5, sendo o maior desvio padrão observado na região Norte ( 18,4 ), e o menor, na região Centro-Oeste ( 14,7 ).
Os demais desvios foram : 14,8 na região Nordeste, 16,8 na região Sudeste e 15,5 na região Sul.
A mediana das notas dos estudantes de todo o Brasil foi 46,0, a mesma da região Norte.
A maior mediana ocorreu na região Sul ( 51,0 ), e a menor, na região Centro-Oeste ( 43,3 ).
As demais medianas foram : 49,3 na região Nordeste ; e 44,5 na região Sudeste.
A nota máxima do Brasil, como um todo, foi 88,9, sendo obtida por, pelo menos, um aluno na região Sudeste.
As demais notas máximas foram : 72,3 na região Norte ; 83,3 nas regiões Nordeste e Sul ; e 80,3 na região Centro-Oeste.
A menor nota mínima foi obtida na região Sudeste ( 4,0 ), enquanto a maior nota mínima foi obtida na região Nordeste ( 22,2 ).
AS demais notas mínimas foram : 10,1 na região Norte ; 15,7 na região Sul ; e 11,6 na região Centro-Oeste.
Observa -se que não existe diferença estatisticamente significativa da menor média das notas do Componente de Conhecimento Específico da região Centro-Oeste ( 43,5 ) em relação à maior média, obtida na região Nordeste ( 49,7 ).
O intervalo de confiança para a média da região Norte ( 21,8 ) é significativamente maior do que de todas as outras.
Assim como os Gráficos 3.1 e 3.2, o Gráfico 3.3, apresentado a seguir, permite uma avaliação do desempenho de concluintes de Tecnologia em Automação Industrial em relação ao Componente de Conhecimento Específico com um histograma da distribuição das notas correspondentes.
Dentre as três distribuições apresentadas, esta é a mais concentrada nas notas baixas.
Esta também é uma distribuição unimodal, com moda no intervalo ( 40 ; 50 ], o mesmo da prova como um todo e abaixo do grupo modal para a Formação Geral.
A Tabela 3.6 apresenta uma comparação dos resultados em relação à Categoria Administrativa e à Organização Acadêmica, agora levando em conta o desempenho de alunos do Componente de Conhecimento Específico da prova.
No que se refere à Organização Acadêmica, a maior média foi das Universidades ( 47,0 ), vindo a seguir a das Faculdades ( 46,1 ) e, depois, a dos Centros Universitários ( 42,0 ).
O maior desvio padrão, e acima do valor para o Brasil, como um todo, foi o das Faculdades ( 16,8 ).
As Faculdades também obtiveram a maior nota máxima ( 88,9 ) e as Universidades, obtiveram maior mediana ( 47,5 ).
As demais notas máximas foram 83,3 para as Universidades e 79,1 para os Centros Universitários, enquanto as demais medianas foram 40,0 para os Centros Universitários e 46,0 para as Faculdades.
As notas mínimas foram : 4,0 para as Universidades ; 7,6 para os Centros Universitários ; e 4,1 para as Faculdades.
Quanto à Categoria Administrativa, observa -se um comportamento semelhante àquele da parte de Formação Geral e à prova como um todo, ou seja, existe diferença 3.2 ANÁLISE DAS QUESTÕES OBJETIVAS Esta seção apresenta estatísticas para o conjunto de questões objetivas e uma análise gráfica de questões selecionadas dos componentes de Formação Geral ( 3.2.1 ) e de Conhecimento Específico ( 3.2.2 ).
São também apresentadas e comparadas as médias das subpopulações caracterizadas por Grande Região. 3.2.1 Componente de Formação Geral A Tabela 3.7 apresenta as estatísticas básicas relativas às oito questões objetivas do componente da prova que abrange a Formação Geral dos estudantes.
A média do Brasil foi 56,9.
A menor média foi encontrada na região Centro-Oeste ( 53,0 ), e a maior, na região Norte ( 64,8 ).
As demais médias foram : 63,3 na região Nordeste ; 56,5 na região Sudeste ; e 57,9 na região Sul.
O desvio padrão do Brasil foi 20,6, sendo o maior desvio padrão encontrado na região Centro-Oeste ( 22,3 ), e o menor, na região Norte ( 14,6 ).
Os demais desvios foram : 18,2 na região Nordeste ; 20,8 na região Sudeste ; e 19,9 na região Sul.
As medianas ( 62,5 ) foram iguais para todas as regiões.
As notas máximas foi 100,0 para todas as regiões com exceção da região Norte onde a nota máxima foi 87,5.
A nota mínimas foi diferente de zero nas Norte ( 37,5 ) e Nordeste ( 12,5 ). estatisticamente significativa entre as médias das IES Públicas ( 53,4 ) e IES Privadas ( 42,5 ).
Mais uma vez, a maior média foi obtida por alunos de IES Públicas de ensino.
Observa -se que existem diferenças estatisticamente significativas ao nível de 95 \% no Componente de Conhecimento Específico entre as notas dos Centros Universitários em relação às Universidades e às Faculdades. 3.2 ANÁLISE DAS QUESTÕES OBJETIVAS Esta seção apresenta estatísticas para o conjunto de questões objetivas e uma análise gráfica de questões selecionadas dos componentes de Formação Geral ( 3.2.1 ) e de Conhecimento Específico ( 3.2.2 ).
São também apresentadas e comparadas as médias das subpopulações caracterizadas por Grande Região. 3.2.1 Componente de Formação Geral A Tabela 3.7 apresenta as estatísticas básicas relativas às oito questões objetivas do componente da prova que abrange a Formação Geral dos estudantes.
A média do Brasil foi 56,9.
A menor média foi encontrada na região Centro-Oeste ( 53,0 ), e a maior, na região Norte ( 64,8 ).
As demais médias foram : 63,3 na região Nordeste ; 56,5 na região Sudeste ; e 57,9 na região Sul.
O desvio padrão do Brasil foi 20,6, sendo o maior desvio padrão encontrado na região Centro-Oeste ( 22,3 ), e o menor, na região Norte ( 14,6 ).
Os demais desvios foram : 18,2 na região Nordeste ; 20,8 na região Sudeste ; e 19,9 na região Sul.
As medianas ( 62,5 ) foram iguais para todas as regiões.
As notas máximas foi 100,0 para todas as regiões com exceção da região Norte onde a nota máxima foi 87,5.
A nota mínimas foi diferente de zero nas Norte ( 37,5 ) e Nordeste ( 12,5 ).
A Tabela 3.8 apresenta o índice de facilidade e o índice de discriminação ( ponto bisserial ) para \textbf{cada} uma das questões objetivas do Componente de Formação Geral.
Quanto ao índice de facilidade, foram usadas as seguintes cores para diferenciar o nível de dificuldade da questão :  Azul para as questões classificadas com índice muito fácil ( \&gt ; =0,86 ), verde para as questões classificadas com índice fácil ( 0,61 a 0,85 ), amarelo para as questões classificadas com médio ( 0,41 a 0,60 ), vermelho para as questões classificadas com difícil ( 0,16 a 0,40 ) e roxo para as questões classificadas com muito difícil ( \&lt ; =0,15 ).
Já quanto ao índice de discriminação, foram usadas as seguintes cores para qualificar a questão :  As questões classificadas com índice fraco receberam a cor vermelho ( \&lt ; =0,19 ), as classificadas com médio receberam a cor amarelo ( 0,20 a 0,29 ), as classificadas com bom receberam a cor verde ( 0,30 a 0,39 ) e as classificadas com muito bom ( \&gt ; =0,40 ) receberam a cor azul.
As questões objetivas do Componente de Formação Geral, segundo o índice de facilidade, foram assim avaliadas : das oito questões, nenhuma teve o índice de facilidade classificado como muito fácil.
Três questões foram tidas como fácil, por terem índice de acertos situado na faixa entre 0,61 e 0,85 ( de 61,0 \% a 85,0 \% de acertos ).
Quatro questões tiveram índice de facilidade médio, situando -se no intervalo entre 0,41 e 0,60, ou seja, houve entre 41,0 \% e 60,0 \% de acertos, enquanto uma questão foi classificada na categoria difícil, situando -se no intervalo entre 0,16 e 0,40.
Por fim, nenhuma das questões apresentou menos de 15 \% de acertos, razão pela qual seria classificada como muito difícil.
Como já comentado, para análise das questões objetivas, relativas à Formação Geral, segundo o poder de discriminação, utilizou -se o índice de discriminação ( ponto bisserial ).
Nesta análise, as questões foram assim avaliadas : seis das oito questões apresentaram índice acima ou igual a 0,40 e, assim, foram classificadas com índice muito bom para esse grupo de estudantes.
Duas questões tiveram índice de discriminação bom, com valor entre 0,30 e 0,39, para esse grupo de estudantes.
Nenhuma questão teve nível médio ou fraco de discriminação para esse grupo de estudantes.
O índice de facilidade variou de 0,36 a 0,79, e o de discriminação, de 0,34 a 0,51.
As seis questões com índice de discriminação muito bom figuraram entre dois níveis de dificuldade desse conjunto : três classificadas na categoria fácil ( questões 3, 4 e 8 ) do índice de facilidade, outras três, na categoria médio ( questões 1, 6 e 7 ).
Em particular, a questão 6 foi a que apresentou o maior poder discriminatório, com índice 0,51, apesar de ser considerada média em termos de facilidade, com uma proporção de 0,54 acertos.
O máximo de acertos foi alcançado pela questão 4 com um índice de facilidade de 0,79.
A questão de número 1 apresentou um índice de facilidade de 0,55, ou seja, pouco mais da metade dos estudantes conseguiu resolvê -la, dentro do universo de participantes.
Seu índice de discriminação foi muito bom ( 0,42 ).
Já a questão 5 obteve índice de discriminação bom, 0,38, porém seu índice de facilidade foi difícil ( 0,36 ).
O Gráfico 3.4, para exemplificar, analisa o comportamento da questão de número 7 de Formação Geral.
Trata -se de uma questão considerada média, com relação a facilidade e a que obteve o segundo maior índice de discriminação dessa parte da prova.
Neste gráfico, \textbf{cada} uma das cinco curvas representa o percentual de respostas em determinada alternativa da questão, em função do número de acertos dos bisserial ).
Nesta análise, as questões foram assim avaliadas : seis das oito questões apresentaram índice acima ou igual a 0,40 e, assim, foram classificadas com índice muito bom para esse grupo de estudantes.
Duas questões tiveram índice de discriminação bom, com valor entre 0,30 e 0,39, para esse grupo de estudantes.
Nenhuma questão teve nível médio ou fraco de discriminação para esse grupo de estudantes.
O índice de facilidade variou de 0,36 a 0,79, e o de discriminação, de 0,34 a 0,51.
As seis questões com índice de discriminação muito bom figuraram entre dois níveis de dificuldade desse conjunto : três classificadas na categoria fácil ( questões 3, 4 e 8 ) do índice de facilidade, outras três, na categoria médio ( questões 1, 6 e 7 ).
Em particular, a questão 6 foi a que apresentou o maior poder discriminatório, com índice 0,51, apesar de ser considerada média em termos de facilidade, com uma proporção de 0,54 acertos.
O máximo de acertos foi alcançado pela questão 4 com um índice de facilidade de 0,79.
A questão de número 1 apresentou um índice de facilidade de 0,55, ou seja, pouco mais da metade dos estudantes conseguiu resolvê -la, dentro do universo de participantes.
Seu índice de discriminação foi muito bom ( 0,42 ).
Já a questão 5 obteve índice de discriminação bom, 0,38, porém seu índice de facilidade foi difícil ( 0,36 ).
O Gráfico 3.4, para exemplificar, analisa o comportamento da questão de número 7 de Formação Geral.
Trata -se de uma questão considerada média, com relação a facilidade e a que obteve o segundo maior índice de discriminação dessa parte da prova. estudantes nessa parte da prova ( Formação Geral/Múltipla Escolha ), antes de possíveis eliminações pelo critério do ponto bisserial.
A curva em vermelho corresponde à alternativa E, a correta para esta questão.
Assim, observa -se que entre os estudantes com menor número de acertos, nessa parte do exame, a situação mais frequente foi a escolha de uma das alternativas incorretas : a alternativa C ( em preto ) ou D ( em roxo ).
Na medida em que o número de acertos aumenta, indicando desempenho melhor nessa parte da prova, aumenta concomitantemente a proporção de estudantes que selecionaram a alternativa correta E, atingindo 100 \% apenas para os estudantes com 8 acertos.
Essa análise permite verificar como a questão discriminou os grupos de desempenho, justificando o alto índice obtido na questão.
Os gráficos relativos às demais questões de Formação Geral constam do Anexo I. 3.2.2 Componente de Conhecimento Específico A Tabela 3.9 apresenta as estatísticas básicas em relação às questões objetivas do Componente de Conhecimento Específico da prova de Tecnologia em Automação Industrial por Grande Região.
A média do Brasil deste componente foi de A média do Brasil deste componente foi de 48,3.
A menor média foi observada na região Norte ( 45,5 ), e a maior, na região Nordeste ( 52,4 ).
O desvio padrão de todo o Brasil foi 17,6, sendo o menor desvio padrão encontrado nas regiões Nordeste e Centro-Oeste ( 15,8 ), e o maior, na região Norte ( 19,5 ).
A mediana de todo o Brasil foi 47,6, a mesma encontrada nas regiões Norte e Sudeste.
Nas demais regiões a mediana foi : 52,4 nas regiões Nordeste e Sul ; e 42,9 na região Centro-Oeste.
A nota máxima da prova foi 90,5, obtida nas questões objetivas do Componente de Conhecimento Específico, por, pelo menos, um aluno nas regiões Sudeste e Sul, e nas demais regiões, a nota máxima da prova foi : 76,2 na região Norte ; e 85,7 nas regiões Nordeste e Centro-Oeste.
A menor nota mínima foi obtida na região Sudeste ( 0,0 ), e nas demais regiões a nota mínima foi : 9,5 nas regiões Norte e Centro-Oeste ; 23,8 na região Nordeste ; e 14,3 na região Sul.
A Tabela 3.10 apresenta os índices de facilidade e discriminação ( ponto bisserial ) das questões objetivas do Componente de Conhecimento Específico para os estudantes de Tecnologia em Automação Industrial.
Para facilitar a diferenciação das questões, usaremos as mesmas cores da Tabela 3.8 para as diferentes classificações dos índices de facilidade e de discriminação.
Dentre as questões objetivas da parte da prova relativa ao Componente de Conhecimento Específico, uma das questões foi anulada pela Comissão, a de número 25.
Desse modo, a classificação, quanto ao índice de facilidade, foi estabelecida com base em 26 questões.
A partir dos índices obtidos, pode -se concluir que um pouco mais de um terço das questões objetivas da prova foi considerado, pelo menos, difícil : das 26 questões válidas, 11 foram classificadas como difícil ( a classificação modal ) ou como muito difícil.
Sete questões foram classificadas como fácil, e outras oito consideradas como médio.
Já quanto aos índices de discriminação das questões objetivas do Componente de Conhecimento Específico da prova, tem -se como resultado a seguinte classificação : dez das 26 questões foram consideradas como boas, enquanto cinco delas tiveram índice de discriminação muito bom.
Assim, para um pouco mais da metade das questões – 15 em 26 – os índices de discriminação foram bom ou muito bom.
Dentre as demais, seis delas foram classificadas como médio e outras cinco como fraco, sendo onze, por conseguinte, a quantidade de questões nos dois patamares mais baixos de discriminação.
Constata -se, assim, que a prova – no que se refere ao Componente de Conhecimento Específico – possuía boa capacidade de discriminar entre aqueles que dominam ou não o conteúdo.
Dentre as questões que alcançaram os maiores índices de discriminação, cinco delas, as de números 11, 12, 20, 27 e 35, foram classificadas com índice de discriminação muito bom, situando -se no intervalo de 0,40 a 0,51 do índice.
Quatro delas ( questões 11, 12, 20 e 27 ) foram classificadas na categoria fácil, quanto ao índice de facilidade, e a outra ( questão 35 ), na categoria médio.
A questão de número 19 foi a mais difícil dentre as 26 questões específicas válidas, com baixo índice de facilidade, apenas 7,0 \% de acertos.
Essa questão apresentou poder discriminatório igualmente baixo, 0,06, o que comprova ter sido a questão mais difícil para os estudantes.
Destaca -se, também, a questão 24, com índice de facilidade 0,21, o que, em termos percentuais, corresponde a 21,0 \% de estudantes que responderam acertadamente, e seu índice de discriminação foi 0,17.
Tais questões foram, portanto, pelo critério ponto bisserial, consideradas inadequadas.
Por isso, as questões 19 e 24 foram eliminadas do cômputo da nota final.
Além destas duas, as demais questões com índice de discriminação fraco, questões 14, 23 e 32 também não foram utilizadas no cômputo final das notas, num total de cinco questões eliminadas.
A título de exemplo das análises do comportamento das questões objetivas, o Gráfico 3.5 analisa a questão 27 do Componente de Conhecimento Específico.
Esta questão foi classificada como fácil.
Apresentou índice de facilidade 0,62, ou seja, 62,0 \% dos estudantes assinalaram acertadamente a opção A, correspondente ao gabarito.
Seu índice de discriminação foi igual a 0,51, classificado como muito bom, também sendo esta questão a que apresentou o maior índice discriminatório.
Neste gráfico, \textbf{cada} uma das cinco curvas representa o percentual de respostas em determinada alternativa da questão 27, em função do número de acertos dos estudantes nessa parte da prova, antes de possíveis eliminações de questões pelo critério do ponto bisserial.
A alternativa correta A, representada no gráfico pela curva em azul, foi escolhida em maiores proporções pelos alunos com desempenho melhor nessa parte da prova.
Já as alternativas incorretas, também denominadas distratores, foram selecionadas, principalmente, por aqueles com notas mais baixas.
Neste caso também a soma não é sempre 100 \% por causa das questões não respondidas ou com mais de uma opção marcada.
Aqueles com nota zero, na sua quase totalidade deixaram esta questão em branco ou marcaram mais de uma alternativa, comportamento considerado inválido.
A proporção de alunos que selecionou a resposta correta E aumenta gradativamente, chegando a atingir 100 \% para 19 e 22 acertos, enquanto a proporção dos que escolheram alternativas incorretas decai, a partir do primeiro acerto, como função do número de acertos nessa parte da prova.
Os gráficos relativos às demais questões do Conhecimento Específico constam do Anexo I. 3.3 ANÁLISE DAS QUESTÕES DISCURSIVAS Esta seção apresenta estatísticas selecionadas e histogramas das Questões Discursivas de Formação Geral ( 3.3.1 ) e Conhecimento Específico ( 3.3.2 ).
São também apresentadas e comparadas as médias de subpopulações caracterizadas por Grande Região. 3.3.1 Componente de Formação Geral As análises dos resultados de desempenho dos estudantes de Tecnologia em Automação Industrial nas duas questões discursivas relativas à Formação Geral, considerando -se as notas de conteúdo e de desempenho linguístico, encontram -se na Tabela 3.11 e no Gráfico 3.6.
Na Tabela 3.11, observa -se que a nota média nesse conjunto de questões foi próxima da obtida nas objetivas.
Os estudantes, de todo o Brasil, obtiveram, em Formação Geral, média 56,9 nas questões objetivas e 49,2 nas questões discursivas.
No entanto, pode -se notar um aumento do desvio padrão de 20,6 nas questões objetivas do Componente de Formação Geral dos alunos de todo o Brasil, para 24,5 nas questões discursivas do mesmo componente.
A maior média foi obtida na região Nordeste ( 54,5 ), e a menor, na região Centro-Oeste ( 37,9 ).
A mediana de todo o Brasil, neste componente, foi 54,0, sendo maior na região Nordeste ( 57,3 ) e nas regiões Sudeste e Sul ( 55,0 em ambas ).
Nas demais regiões foi menor que a mediana nacional : Norte ( 52,5 ) e Centro-Oeste ( 43,5 ).
A nota máxima ( 97,0 ) foi obtida por, pelo menos, um aluno na região Sudeste ; sendo 88,0 na região Norte, 94,0 na região Nordeste, 95,0 na região Sul e 96,0 na região Centro-Oeste.
A nota mínima ( 0,0 ) foi a mesma em todas as regiões do Brasil, sem exceção.
O Gráfico 3.6 representa a distribuição das notas nas questões discursivas do Componente de Formação Geral.
A moda desta distribuição ocorre no intervalo ( 60 ; 70 ].
Destacam -se, também, o intervalo ( 50 ; 60 ], com distribuição próxima ao valor modal, e o intervalo [ 0 ; 10 ] com distribuição próxima a 12,5 \% do total de notas, sendo que no intervalo [ 0 ; 10 ] incluem -se, além da nota zero, a frequência de alunos que deixaram as questões discursivas de Formação Geral em branco.
Na sequência, os resultados verificados para \textbf{cada} uma das questões discursivas de Formação Geral serão apresentados, estabelecendo -se relações com os conteúdos abordados em \textbf{cada} uma delas.
Os comentários da Banca de docentes corretores, a respeito do observado na correção das respostas dos estudantes, suas impressões e conclusões serão apresentados junto à análise de \textbf{cada} questão.
Cumpre esclarecer que, tendo em vista que as questões discursivas de Formação Geral são padronizadas, ou seja, constam de todas as provas, os comentários da Banca são os mesmos para todas as carreiras acadêmicas, sendo direcionados a todos os estudantes que participaram do ENADE/2014.
A seguir, serão analisados os desempenhos da Área de Tecnologia em Automação Industrial nas duas questões discursivas de Formação Geral do 3.3.1.1 Análise de Conteúdo da Questão Discursiva 1 do Componente de Formação Geral Os dados de Tecnologia em Automação Industrial, obtidos a partir das respostas à questão 1, encontram -se na Tabela 3.12 e no Gráfico 3.7.
Nessa questão – de melhor desempenho dentre as duas de Formação Geral – os alunos, de todo o Brasil, tiveram média 51,0.
A maior média para a questão 1 foi obtida na região Nordeste ( 56,6 ), e a menor, na região Centro-Oeste ( 39,7 ).
Quanto à variabilidade das notas, o desvio padrão de todo o Brasil foi 31,0.
O menor desvio padrão foi obtido na região Nordeste ( 25,9 ), e o maior na região Norte ( 34,4 ).
As medianas das regiões Sudeste, Sul e do Brasil, como um todo, foram iguais ( 55,0 ).
Nas demais regiões a mediana foi : 65,0 na região Norte ; 60,0 na região Nordeste ; e 45,0 na região Centro-Oeste.
As notas máximas e mínimas da questão discursiva 1 foram as mesmas para todas as regiões do Brasil, respectivamente, 100,0 e 0,0, sem exceção.
O Gráfico 3.7 mostra a distribuição das notas na questão discursiva 1 do Componente de Formação Geral.
Observa -se que a maior frequência corresponde aos alunos que obtiveram nota no intervalo ( 70 ; 80 ], correspondendo à moda da distribuição.
Destacam -se, também, os alunos que deixaram a questão em branco e o intervalo ( 40 ; 50 ], com mais de 12,5 \% de frequência em \textbf{cada} uma das situações. 3.3.1.2 Comentários sobre a correção de Conteúdo das respostas à Questão Discursiva 1 O enunciado apresentado da Questão 1 era claro e foi redigido em linguagem acessível.
A questão abordava o tema da mobilidade urbana, relacionando -a aos ideais de desenvolvimento urbano sustentável estabelecidos pela ONU, do ponto de vista ecológico, cultural, político, institucional, social e econômico.
O enunciado apresentava uma tabela, sem referência à fonte de publicação, que detalhava percentualmente as modalidades de deslocamento da população urbana brasileira, com base em duas distinções : motorizado versus não motorizado ; coletivo versus individual.
O comando da questão solicitava, então, que o estudante desenvolvesse um texto dissertativo composto de duas partes : análise das consequências do transporte motorizado para o desenvolvimento sustentável e elaboração de duas propostas de intervenção pautadas pelo incentivo ao uso da bicicleta, meio de transporte não motorizado e, portanto, não poluente.
A temática dos efeitos da industrialização, com foco na produção de automóveis, e consequências para o meio ambiente é atual.
Além disso, a questão provoca a reflexão sobre a necessidade de adequação das cidades ao número de habitantes, no contexto da utilização dos meios de transportes de massa : metroviário, ferroviário, por ônibus.
Esperava -se, assim, que o concluinte recorresse a diferentes argumentos para analisar as consequências do transporte motorizado, explicitados pelo texto motivador, e propusesse duas ações de intervenção, as quais considerasse relevantes para o incremento ao uso de bicicleta, que, preferencialmente, deveriam manter coerência com a argumentação desenvolvida na primeira parte da resposta.
O conflito entre a busca pelo transporte individual e o coletivo é inerente a todas as grandes cidades do mundo.
No Brasil, esse conflito emerge com maior gravidade, como efeito do aumento da quantidade de veículos automotores ( automóveis e motocicletas ) e da baixa qualidade dos transportes coletivos.
Tal associação incentiva o uso de meios de transporte individual em detrimento da utilização dos coletivos.
O padrão de respostas era tecnicamente adequado ao enunciado e amplo o suficiente para abranger as variações de respostas apresentadas.
Deve ser ressaltado que a proteção ao ambiente ecologicamente equilibrado é tema conhecido dos concluintes dos cursos universitários, que deveriam possuir informações suficientes para desenvolver uma resposta abordando pelo menos dois dos aspectos listados no padrão de respostas como possibilidades de consequências ( item ‘ a ’ ) e duas propostas de ação de intervenção aceitáveis ( item ‘ b ’ ).
Segundo o padrão, concluintes poderiam desenvolver, quanto ao item ‘ a ’, os seguintes pontos : aumento da emissão de poluentes atmosféricos ; aumento da emissão de gases de efeito estufa ( CO2- dióxido de carbono, CO - monóxido de carbono, O3 - Ozônio ) ; aumento da poluição visual e sonora ; aumento da temperatura local e global ; aumento do consumo de combustíveis ; aumento de problemas de saúde ( cardíacos, respiratórios, dermatológicos ) ; aumento da frota de veículos, promovendo congestionamentos urbanos ; diminuição de áreas verdes ; desmatamento ; aumento das áreas de impermeabilização, resultando em enchentes, diminuição da infiltração da água e recarga de lençóis freáticos ; elevação dos custos de manutenção das cidades ( metroferrovias, rodovias, tratamento de água, limpeza da cidade, etc. ) ; necessidade de ampliação de vias trafegáveis ; necessidade de ampliação de áreas de estacionamento.
Em relação ao item ‘ b ’, os estudantes poderiam considerar os seguintes aspectos : construção de vias exclusivas para bicicletas ( ciclovias e ciclofaixas ) ; proposição de formas de integração entre o transporte por bicicletas, o metroviário e os ônibus coletivos, a fim de garantir segurança e conforto em momentos de adversidades climáticas e relevo \textbf{acidentado} ; pontos de aluguel e/ou empréstimos de bicicletas ; construção de bicicletários ; investimentos na segurança pública ; políticas de incentivo ao uso de bicicletas ( educação ambiental, qualidade de vida, saúde, propaganda ) ; implementação de políticas de crédito e de redução de custos das bicicletas.
Para o item ‘ a ’, foi encontrado um número expressivo de respostas adequadas.
As consequências mais citadas \textbf{estavam} relacionadas com a emissão de poluentes e o aumento dos engarrafamentos.
Houve críticas frequentes à qualidade do transporte público e à necessidade de locomoção rápida diante da distância das residências aos locais de trabalho.
Em tom de reivindicação, muitos estudantes defenderam solução para os problemas do transporte coletivo, quer metroviário, ferroviário ou por ônibus.
O uso da bicicleta também foi considerado por muitos como um meio de transporte benéfico à saúde, mas alguns apontaram a impossibilidade de sua utilização pelos idosos e deficientes.
No concernente ao item ‘ b ’, de modo quase uniforme citou -se a necessidade de ciclovias e ciclofaixas, bem como de se garantir segurança para os \textbf{ciclistas}.
A necessidade de manutenção das ciclovias e ciclofaixas existentes em algumas cidades foram também \textbf{lembradas}.
Com menos frequência também houve a indicação de maior conscientização da população, inclusive pela mídia, e pelo incentivo à aquisição desses veículos não motorizados, por meio da diminuição de impostos para redução do custo de aquisição.
Houve quem sugerisse, inclusive, uma política de doação de bicicletas.
Alguns mencionaram a necessidade de se assegurar integração das ciclovias com os demais modais, com a construção de bicicletários seguros diante das grandes distâncias a percorrer, notadamente nas metrópoles.
Também com base na dificuldade de o veículo não motorizado vencer as longas distâncias que o trabalhador deve percorrer nos centros urbanos para chegar ao local onde exerce a sua atividade produtiva, alguns alunos argumentaram contra a utilização da bicicleta como forma de resolução da poluição ambiental.
O principal equívoco cometido na sugestão de ações ( item ‘ b ’ ) foi o de mencionar a necessidade de melhoria do meio ambiente dissociado do uso de bicicletas, como se requeria no comando da questão.
Nesse sentido, foram encontradas respostas que indicavam a caminhada como opção.
Foram consideradas respostas fracas ( notas de zero a 30 ) aquelas meramente opinativas, ou que apenas transcreveram elementos do enunciado, sem qualquer acréscimo ou apreciação crítica.
As respostas medianas ( notas entre 35 e 70 ) foram as que não apresentavam algum desenvolvimento para os tópicos citados ( duas consequências e duas ações de intervenção ) ou por falta de alguns deles.
Já as respostas boas ( notas de 75 a 100 ) continham os tópicos considerados pelo padrão de respostas, com algum desenvolvimento, pelo menos, regular de argumentação.
As respostas dos estudantes revelaram um adequado conhecimento quanto aos efeitos do transporte motorizado para o ambiente e que a temática da ecologia parece ser tema bem compreendido, aparecendo nas respostas como uma preocupação para as futuras gerações.
Nesse sentido, algumas políticas públicas foram mencionadas apropriadamente para assegurar o desenvolvimento sustentável pelos respondentes, tais como : educação, habitação, saneamento, saúde, transporte e segurança.
No entanto, de modo geral, as respostas demostraram que os estudantes têm muita dificuldade na expressão escrita do pensamento, como se pode constatar pela avaliação do desempenho linguístico que ficou a cargo de uma banca específica, formada por profissionais da Área de Língua Portuguesa. 3.3.1.3 Análise de Conteúdo da Questão Discursiva 2 do Componente de Formação Geral A Tabela 3.13 mostra que o desempenho médio dos estudantes na questão discursiva 2 ( média 46,6 ) foi inferior ao obtido na questão discursiva 1 ( média 51,0 ).
A região Nordeste foi aquela cuja média, nessa questão, foi maior ( 50,5 ), e a de menor média foi a região Centro-Oeste ( 34,9 ).
O desvio padrão de todo o Brasil foi 30,1, inferior ao obtido na questão discursiva 1 ( 31,0 ).
O maior desvio nessa questão foi obtido na região Norte ( 36,4 ), enquanto o menor foi obtido na região Sudeste ( 29,5 ).
A mediana de todo o Brasil foi 50,0, a mesma das regiões Nordeste, Sudeste e Sul, enquanto nas regiões Norte e Centro-Oeste a mediana foi, respectivamente, 45,0 e 30,0.
As notas máximas ( 100,0 ) foram as mesmas em todas as regiões do Brasil, assim como as notas mínimas ( 0,0 ), sem exceção. bicicletas, como se requeria no comando da questão.
Nesse sentido, foram encontradas respostas que indicavam a caminhada como opção.
Foram consideradas respostas fracas ( notas de zero a 30 ) aquelas meramente opinativas, ou que apenas transcreveram elementos do enunciado, sem qualquer acréscimo ou apreciação crítica.
As respostas medianas ( notas entre 35 e 70 ) foram as que não apresentavam algum desenvolvimento para os tópicos citados ( duas consequências e duas ações de intervenção ) ou por falta de alguns deles.
Já as respostas boas ( notas de 75 a 100 ) continham os tópicos considerados pelo padrão de respostas, com algum desenvolvimento, pelo menos, regular de argumentação.
As respostas dos estudantes revelaram um adequado conhecimento quanto aos efeitos do transporte motorizado para o ambiente e que a temática da ecologia parece ser tema bem compreendido, aparecendo nas respostas como uma preocupação para as futuras gerações.
Nesse sentido, algumas políticas públicas foram mencionadas apropriadamente para assegurar o desenvolvimento sustentável pelos respondentes, tais como : educação, habitação, saneamento, saúde, transporte e segurança.
No entanto, de modo geral, as respostas demostraram que os estudantes têm muita dificuldade na expressão escrita do pensamento, como se pode constatar pela avaliação do desempenho linguístico que ficou a cargo de uma banca específica, formada por profissionais da Área de Língua Portuguesa. 3.3.1.3 Análise de Conteúdo da Questão Discursiva 2 do Componente de Formação Geral A Tabela 3.13 mostra que o desempenho médio dos estudantes na questão discursiva 2 ( média 46,6 ) foi inferior ao obtido na questão discursiva 1 ( média 51,0 ).
A região Nordeste foi aquela cuja média, nessa questão, foi maior ( 50,5 ), e a de menor média foi a região Centro-Oeste ( 34,9 ).
O desvio padrão de todo o Brasil foi 30,1, inferior ao obtido na questão discursiva 1 ( 31,0 ).
O maior desvio nessa questão foi obtido na região Norte ( 36,4 ), enquanto o menor foi obtido na região Sudeste ( 29,5 ).
A mediana de todo o Brasil foi 50,0, a mesma das regiões Nordeste, Sudeste e Sul, enquanto nas regiões Norte e Centro-Oeste a mediana foi, respectivamente, 45,0 e 30,0.
As notas máximas ( 100,0 ) foram as mesmas em todas as regiões do Brasil, assim como as notas mínimas ( 0,0 ), sem exceção.
O Gráfico 3.8 mostra a distribuição das notas na questão discursiva 2 do Componente de Formação Geral.
Observa -se que a maior frequência corresponde aos alunos que obtiveram nota situada no intervalo ( 40 ; 50 ].
Destacam -se também, os alunos que deixaram a questão 2 em branco ( acima de 15,0 \% ). 3.3.1.4 Comentários sobre a correção de Conteúdo das respostas à Questão Discursiva 2 A questão 2 abordava o tema da violência urbana a partir de um texto motivador recolhido em um site de um jornal brasileiro.
O fragmento relatava a prisão de três jovens por agressão a um outro jovem, menor de idade, guardador de carros.
A notícia trazia informações sobre o estado do jovem agredido e sobre as possíveis penas a serem aplicadas aos agressores.
O enunciado solicitava, então, que o estudante desenvolvesse um texto dissertativo composto de duas partes : análise de duas causas do tipo de violência descrito no texto e apresentação de dois fatores que contribuiriam para evitar a agressão relatada.
Como o texto motivador reduzia -se ao relato de um episódio de violência urbana, o enunciado não oferecia argumentos que ajudassem o estudante a elaborar uma análise das possíveis causas do fato.
As propostas a serem apresentadas, no sentido de se evitar em situações como a relatada, deveriam manter coerência com a análise das causas, desenvolvida na primeira parte da resposta.
O enunciado era claro e a matéria permitia identificar tratar -se de um conflito tipicamente urbano no qual se envolveram jovens de baixa renda em disputa por um território onde exerciam trabalho informal.
Sendo essa uma das hipóteses albergada sobre o amplo tema da violência urbana no qual se desenvolve o relato.
A temática da questão é bastante conhecida, não somente por debates no âmbito universitário, especialmente na Área de Ciências Sociais, bem como pela ênfase que é dada pela mídia onde assume contornos, muitas vezes, maiores do que o problema em si.
Assim, por tratar -se de tema bastante presente e debatido no cotidiano das pessoas, a questão pode ser considerada de baixa complexidade.
O padrão de resposta apresentava uma gama de possibilidades de causas da violência aceitas como corretas para o item ‘ a ’, tais como : problemas relacionados à educação ( baixa escolaridade, evasão escolar, qualidade da educação, distanciamento entre escola e a realidade social, tempo de permanência na escola ) ; desigualdades socioculturais ( gênero, etnia, economia, etc. ) ; desemprego e falta de qualificação profissional ; precariedade da segurança pública ; uso de drogas ; \textbf{desvalorização} da vida humana ; \textbf{banalização} da violência ; sensação de impunidade ; ausência de políticas sociais ; degradação da vida urbana ; desconhecimento ou desrespeito aos direitos humanos e constitucionais ; desestruturação familiar ; \textbf{desvalorização} de princípios éticos e morais.
No concernente ao item ‘ b ’, o padrão também previa a possibilidade de apresentação de diversos fatores que podem contribuir para evitar a violência : políticas de segurança mais efetivas ; políticas públicas de melhoria das condições socioeconômicas ; maior consciência cidadã e respeito à vida ; melhor distribuição de renda ; melhoria da educação ( aumento da escolaridade, redução da evasão escolar, qualidade da educação, aproximação entre a escola e a realidade social, aumento do tempo de permanência na escola ) ; aumento da oferta de emprego e melhoria da qualificação profissional ; medidas preventivas ao uso de drogas ; maior eficácia do sistema judiciário ; revisão da legislação penal ; valorização de princípios éticos, morais e familiares.
O enunciado da questão permitia que os estudantes, em suas respostas, pautassem o problema por diversos ângulos.
Dessa forma, o padrão buscou indicar diversas possibilidades a serem identificadas pelo corretor na leitura da resposta do estudante.
As notas foram atribuídas considerando o desenvolvimento da resposta e a indicação de ao menos dois dos itens exigidos pelo comando da questão nos itens ’ a ’ ( duas causas ) e ‘ b ’ ( dois fatores ), dentre os considerados pelo padrão de respostas.
No entanto, o padrão exigia que as respostas \textbf{estivessem} afinadas com o perfil profissional previsto no Art. 3º da Portaria nº 255, de 02 de junho de 2014, onde se lê : Art. 3º No componente de Formação Geral serão considerados os seguintes elementos integrantes do perfil profissional : atitude ética ; comprometimento social ; compreensão de temas que transcendam ao ambiente próprio de sua formação, relevantes para a realidade social ; espírito científico, humanístico e reflexivo ; capacidade de análise crítica e integradora da realidade ; e aptidão para socializar conhecimentos em vários contextos e públicos diferenciados.
Nesse sentido, teses repressivas não conformadas com o texto constitucional e amparados pelos direitos humanos, não foram consideradas.
Questões como a majoração das penas, instituição de penas mais duras e \textbf{criminalização} da conduta dos menores desviados, só foram consideradas no caso de serem tratadas como menção a um debate que \textbf{está} presente, nunca como medida a ser aplicada de forma arbitrária e inconstitucional.
Também não foram previstas visões mais autoritárias e repressoras como a defesa de proibição do trabalho informal, como é o caso dos “ flanelinhas ”, e a defesa de não fornecer dinheiro aos “ guardadores ” de carro, identificado como esmolas por alguns.
Não foram consideradas respostas que meramente reproduziam o enunciado sem qualquer reflexão ou análise.
Aquelas que utilizaram os elementos do enunciado, mesmo que apenas transcritos, mas desenvolveram algum raciocínio pessoal ou indicaram causas e/ou fatores foram valoradas.
As causas apontadas com maior frequência para a violência foram : a ausência de educação e o desemprego, com soluções correlatas de educação profissional e de qualidade e abertura de postos de trabalho.
A questão das drogas também foi bastante mencionada, e a solução correlata mais indicada foi o acolhimento de usuários de drogas.
Outro tema bastante citado foi a dissolução das famílias e dos valores familiares.
Tendo em vista que o texto falava de moradores de rua, a ausência de moradia apareceu em muitas respostas, bem como a disputa territorial e a necessidade de sobrevivência.
Nesse sentido, as propostas apontavam para a necessidade de ajuda da Área de Serviço Social das prefeituras e atendimento psicológico gratuito.
Numa ótica mais repressora, surgiram críticas quanto à leniência das leis, acompanhadas da indicação da necessidade de atuação mais rigorosa do Poder Judiciário, como também da aplicação de penas maiores, negativa de esmolas, proibição do trabalho informal ( guardador de carros ), dentre outras.
Algumas respostas indicaram como causa o modelo capitalista, o consumismo e a ganância.
Outras indicaram o trabalho infantil e a correlata necessidade de aplicação do ECA ( Estatuto da Criança e do Adolescente ).
Em relação ao tema violência, inúmeras respostas apresentaram visões de intolerância quanto aos menores infratores, postulando a diminuição da idade para efeitos de \textbf{criminalização}, bem como defendendo a majoração de penas.
Essa perspectiva revela uma cisão social, de \textbf{certa} forma estimulada pelo tipo de divulgação da mídia quanto aos temas em foco, o que, sem dúvida, ficou evidenciado em algumas respostas.
A existência de inúmeras respostas pugnando pela maior repressão estatal e pelo aumento das penas, acrescido da busca pela diminuição da idade penal, revela que muitos estudantes não foram atingidos por uma formação superior que se exige humanista, ética e comprometida socialmente, que busque os fundamentos para resolver o problema da violência por meio de práticas democráticas e de inclusão social.
As respostas desse tipo também revelaram um alto grau de intolerância para com os moradores de rua, associando -os muitas vezes ao tráfico de drogas, à exploração indevida do espaço público, bem como a crimes como extorsão dos motoristas quando pedem dinheiro para guardar os veículos.
Nessa trilha, verifica -se lacuna na discussão de temas importantes para a formação de profissionais de nível superior tais como : sociodiversidade, multiculturalismo e violência ; tolerância/intolerância ; inclusão/exclusão.
Para aqueles que responderam em conformidade com o padrão de resposta verificou -se uma plena formação social, conclamando pelo auxílio não só do Estado, mas também da Sociedade no amparo daqueles que não têm as mínimas condições de sobrevivência e que \textbf{estão} nas ruas. 3.3.1.5 Análise de Língua Portuguesa das Questões Discursivas do Componente de Formação Geral Os dados de Tecnologia em Automação Industrial, obtidos a partir das respostas às questões discursivas do Componente de Formação Geral, no que tange à Língua Portuguesa, encontram -se na Tabela 3.14 e no Gráfico 3.9.
Nesse aspecto, os alunos, de todo o Brasil, tiveram média 51,0.
A maior média com respeito à Língua Portuguesa foi obtida na região Nordeste ( 58,3 ), e a menor, na região Centro-Oeste ( 40,1 ).
Quanto à variabilidade das notas, o desvio padrão de todo o Brasil foi 23,0.
O menor desvio padrão foi obtido na região Nordeste ( 19,2 ) e o maior desvio padrão foi obtido na região Norte ( 28,0 ).
A mediana das notas de Língua Portuguesa foi 57,5 para as regiões Sudeste, Sul e o Brasil com um todo.
As demais medianas foram : 62,5 na região Norte ; 65,0 na região Nordeste ; e 48,8 na região Centro-Oeste.
A nota máxima para todo o Brasil foi de 90,0, com, pelo menos, um aluno tirando essa nota nas regiões Sudeste e Sul.
As demais notas máximas foram : 80,0 na região Norte ; e 85,0 nas regiões Nordeste.
Além disso, a nota mínima foi zero em todas as regiões do país, sem exceção.
O Gráfico 3.9 mostra a distribuição das notas de Língua Portuguesa do Componente de Formação Geral.
Observa -se que a maior frequência corresponde aos alunos que obtiveram nota no intervalo ( 60 ; 70 ]. 3.3.1.6 Comentários sobre a correção das respostas de Formação Geral com respeito à Língua Portuguesa Ao encaminhar as questões 1 e 2 na direção da produção de um texto dissertativo, esperava -se que o estudante utilizasse seus conhecimentos sobre o assunto e estruturasse seus textos de acordo com as exigências do registro formal próprio dessa situação comunicativa.
Essa configuração determina exigências quanto : à adequação da seleção vocabular, ao desenvolvimento do conteúdo, à estruturação sintática dos períodos, à organização lógica das ideias, à utilização de procedimentos de encadeamento textual e referenciação, à obediência às exigências morfossintáticas próprias da modalidade escrita da norma-padrão, ao respeito às regras ortográficas e às regras de acentuação gráfica.
O padrão de respostas utilizado na avaliação considerou os aspectos relevantes ao bom desempenho linguístico como competências distintas, de modo a permitir um mapeamento detalhado do domínio dos recursos disponíveis na Língua Portuguesa para a comunicação escrita formal.
Com base nesse objetivo, foram avaliados os seguintes aspectos : a ) Estruturação textual condizente com o gênero solicitado e o modo de organização textual expositivo adequado ao gênero – essa competência envolve : a estruturação sintática condizente com o padrão da modalidade escrita formal da língua portuguesa, de modo a garantir a clareza necessária ; a distribuição do conteúdo do texto em parágrafos, de modo a garantir a sua organização temática ; a utilização de operadores discursivos que contribuam para a progressão temática do texto, estabelecendo relações lógicas entre as ideias apresentadas, tanto do ponto de vista intrafrasal, como do interfrasal ; a utilização de procedimentos de referenciação lexical e pronominal que permitam a retomada de referentes textuais ; o respeito às regras de pontuação como fator de estruturação do período.
Espera -se, portanto, que o estudante recorra a procedimentos linguístico-discursivos para organizar seu texto, permitindo o encadeamento lógico entre suas partes de forma a garantir a progressão e a coerência textuais.
Isso significa que os seguintes procedimentos foram penalizados, de acordo com o padrão de respostas proposto : • a estrutura lógico-gramatical do texto fica comprometida pela elaboração de frases fragmentadas ; • sequência justaposta de ideias sem encaixamentos sintáticos, reproduzindo hábitos da oralidade ; • elaboração de frase com apenas oração subordinada, sem oração principal ; • emprego equivocado do conector ( preposição, conjunção, pronome relativo, alguns advérbios e locuções adverbiais ) que não estabeleça relação lógica entre dois trechos do texto e prejudique a compreensão da mensagem ; • emprego do pronome relativo sem a preposição, quando obrigatória ; • repetição ou substituição inadequada de palavras sem utilização dos recursos oferecidos pela língua ( pronome, advérbio, artigo, sinônimo ) ; • emprego inadequado dos pronomes relativos “ cujo(a) ” e “ onde ” ; • utilização inadequada dos sinais de pontuação que comprometa a clareza textual. b ) Respeito às convenções ortográficas da norma-padrão da Língua Portuguesa – essa competência envolve o domínio das regras de acentuação gráfica e da grafia padrão das palavras ( com ausência de abreviaturas próprias da linguagem da internet ), de acordo com as convenções estabelecidas pela legislação em vigor e consubstanciadas no Vocabulário Ortográfico da Língua Portuguesa, editado pela Academia Brasileira de Letras ( com aceitação da legislação anterior, no caso das regras relativas ao uso do hífen e da acentuação gráfica ).
Espera -se que o participante :  grafe corretamente as palavras ;  respeite as regras de acentuação gráfica ;  empregue maiúsculas em início de frase, em nomes próprios de pessoas, lugares ou instituições ;  evite abreviações como p/, vc, tb, pq, \textbf{tá}, né, usadas muitas vezes em escrita informal e na internet ;  obedeça às regras de separação de sílabas no final da linha. c ) Domínio dos diferentes aspectos morfossintáticos próprios da modalidade escrita formal da norma-padrão da Língua Portuguesa – essa competência envolve : a concordância nominal, a concordância verbal, a regência nominal, a regência verbal, a flexão nominal, a flexão verbal, a correlação entre os tempos verbais, a colocação pronominal e a utilização de sinais de pontuação que contribuam para a organização lógica da frase e do texto.
Espera -se que o participante :  flexione o verbo para estabelecer concordância de número com o sujeito da frase ;  flexione o artigo, o adjetivo e o pronome para concordar em número e em gênero com o substantivo a que se referem ;  observe a regência nominal e a verbal, utilizando a preposição adequada depois de um substantivo, um verbo ou um adjetivo ;  empregue adequadamente o acento grave indicador de crase nos casos em que se fizer necessário ;  obedeça às regras de colocação pronominal ( próclise e ênclise ), distintas dos hábitos da oralidade ou da escrita informal ; coloquiais, evite repetição desnecessária de palavras e utilize um vocabulário mais formal, adequado ao texto de caráter dissertativo.
A escolha dessas competências para subsidiar o processo de avaliação apoia -se na concepção de que, no desempenho dos graduandos, a modalidade escrita tem apresentado uma intensa simplificação, originada no padrão da modalidade oral da Língua Portuguesa.
No caso do texto de base dissertativa, inscrito em um registro formal, a distância entre as duas modalidades é ainda maior, o que provoca situações de hipercorreção ( desvios provocados pela incorporação indevida de uma regra da norma-padrão ) e de truncamentos sintáticos ( estruturas frasais incompreensíveis devido à complexidade sintática própria da modalidade escrita ).
Observam -se, então, os seguintes aspectos que marcam essa distinção entre as duas modalidades, devido à excessiva simplificação da modalidade falada : a ) redução drástica de estruturas subordinadas, compensada pelo aumento na frequência de estruturas coordenadas e absolutas, por um lado, ou pela elaboração de estruturas truncadas pelo excesso de ideias sem a devida conexão subordinativa ; b ) redução no uso de conectores para expressar relações lógicas essenciais à construção do texto, substituídas pela exigência de inferência por parte do interlocutor para suprir a sua ausência ; c ) redução \textbf{cada} vez maior do uso do subjuntivo, ao lado da ampliação do uso do indicativo combinado a estruturas frasais coordenadas ou absolutas ; d ) empobrecimento do processo de referenciação, com a repetição exaustiva de pronomes ou nomes ; e ) simplificação extrema da marcação da categoria tempo na morfologia verbal ; f ) falta de domínio de vocabulário mais abstrato e de maior complexidade, essencial ao desenvolvimento do processo dissertativo ; g ) redução drástica no emprego da acentuação gráfica, processo intensificado pela divulgação imprecisa das mudanças promovidas pelo último acordo ortográfico.
Os aspectos macroestruturais da elaboração do texto não foram avaliados neste processo, para não penalizar duplamente os estudantes, já que a banca de formação geral, composta por profissionais de diferentes Áreas do conhecimento, encarregou -se da avaliação do conteúdo desenvolvido nas questões.
São eles : progressão temática, coerência na relação com os conhecimentos de mundo dos usuários da língua, inteligibilidade, atendimento ao solicitado no enunciado do ponto de vista do desenvolvimento do conteúdo, entre outros.
A grade de avaliação do desempenho linguístico considerou, portanto, três grandes grupos de competências, segundo os aspectos explicitados anteriormente : 1.
Aspectos ortográficos : domínio das convenções ortográficas : grafia de vogais e consoantes, uso de maiúsculas e minúsculas, emprego do hífen e acentuação gráfica ; 2.
Aspectos textuais : domínio dos procedimentos de estruturação textual do ponto de vista microestrutural : organização interna dos períodos, emprego de conectores para a articulação lógica entre os períodos e entre os parágrafos, emprego de marcas de referenciação lexical e pronominal ; utilização dos sinais de pontuação que contribuem para a organização lógica da frase. 3.
Aspectos morfossintáticos e vocabulares : domínio das regras de caráter morfossintático estabelecidas como modelares do ponto de vista da modalidade escrita formal da norma-padrão da Língua Portuguesa : concordância nominal e verbal, regência nominal e verbal, colocação pronominal, flexão nominal e verbal, correlação entre tempos e modos verbais, ausência de marcas de oralidade.
A seleção vocabular adequada à modalidade escrita formal da Língua Portuguesa foi incorporada a essa última competência, tendo em vista a intersecção entre as duas do ponto de vista das exigências do registro formal da modalidade escrita da norma-padrão.
Os resultados da avaliação correspondem aos seguintes aspectos observados em \textbf{cada} competência : Aspectos ortográficos : A correção foi realizada classificando os textos em cinco níveis, nível zero a quatro.
O desempenho dos estudantes ficou concentrado, majoritariamente, no nível 3 ( de 4 a 7 desvios ), devido, principalmente, aos desvios de acentuação.
O número de desvios de grafia foi reduzido.
Diferentemente do resultado do ENADE/2013, aumentou a porcentagem de textos sem desvios ( enquadrados no nível 4 ) e diminuiu a porcentagem de textos enquadrados no nível 2.
Alguns textos, com pior desempenho, foram enquadrados no nível 1 ( de 8 a 12 desvios ), enquanto o nível zero foi atribuído a pouquíssimos casos, já que ele revela falta absoluta de domínio das convenções ortográficas.
Observou -se, portanto, que existe uma diferença muito grande de desempenho nos dois aspectos analisados : baixo índice de desvios ortográficos e grande índice de desvios de acentuação.
Em vários casos, ocorre ausência completa de acentuação gráfica.
Os resultados revelam, portanto, que a tendência dominante entre os universitários brasileiros é a eliminação da acentuação gráfica, provavelmente motivada pela vivência dos jovens relacionada aos aplicativos de comunicação via internet ( redes sociais e e-mails ).
Nesse tipo de comunicação, devido ao ritmo intenso de troca de mensagens, o uso de acentos gráficos foi praticamente abolido.
Outro fator que pode ter relação com essa tendência é a ausência de esclarecimento dos meios de comunicação, das autoridades e das escolas sobre as decisões do Acordo Ortográfico da Língua Portuguesa de 1990, gerando um estado de indefinição para os estudantes.
Os casos mais sistemáticos de eliminação do acento indicador da sílaba tônica são :  palavras proparoxítonas ( “ ridiculos ”, “ publicas ”, “ lideres ”, “ politicos ” ) ;  palavras paroxítonas terminadas em ditongo crescente ( “ homicidios ”, “ latrocinio ”, “ individuo ”, “ dependencia ” ) ;  palavras oxítonas ( “ ninguem ”, “ esta ”, “ ate ”, “ ai ” ).
Por outro lado, destaca -se o uso indevido do acento gráfico em determinadas palavras, como observado nas grafias * “ jornáis ”, * “ telejornaís ”, * “ propíciar ”, * “ medídas ”, * “ dígnidade ”, * “ cídades ”.
Quanto ao domínio das convenções relativas à grafia das palavras, observam -se desvios como : a hipercorreção pela escolha de “ e ” no lugar de “ i ”, por influência de hábitos da oralidade ( * “ enumeros ” por “ inúmeros ”, * “ entevenção ” por “ intervenção ” ) ; a eliminação do “ r ” marcador do infinitivo verbal ( * “ esta ” no lugar de “ \textbf{estar} ” ).
Outros casos de desvios de grafia relacionados à variação diastrática podem ser observados em *“estrupo ”, *“automovís ”, *“viensse ”, *“camihada ”.
São muito frequentes os seguintes desvios de caráter ortográfico, com repercussão morfossintática :  eliminação da marca de infinitivo ( -r- ) e substituição por acento agudo ( “ \textbf{está} ” no lugar de “ \textbf{estar} ” ) ;  confusão entre “ ão ” e “ am ” nas formas verbais ( “ invadão ” no lugar de “ invadam ” e “ estam ” no lugar de “ \textbf{estão} ” no presente do indicativo ; “ estaram ” no lugar de “ \textbf{estarão} ” no futuro do indicativo ) ;  confusão entre a grafia do verbo “ haver ” ( “ há ” ) e o artigo definido ou a preposição “ a ” ;  uso de hífen para separar pronome átono – tanto uso indevido quanto omissão ( no pretérito imperfeito e futuro do subjuntivo : “ evitar -mos ” no lugar de “ evitarmos ”, “ percebesse ” no lugar de “ percebe -se ” e vice-versa ) ; Observam -se, também, muitos casos de inadequação no uso da maiúscula : ausência de diferença entre a primeira letra e as outras, em início de período, principalmente dos estudantes que adotam a escrita em letra de imprensa ; utilização de maiúscula para destacar determinadas palavras-chave do texto, como “ Violência ”, “ Brasileiros ”, “ Fatores ”, “ Ozônio ”, “ Sustentável ”.
Vale observar, também, que, ao contrário do que se esperava, não apareceram abreviaturas próprias do “ internetês ”, ou seja, dos hábitos de comunicação escrita adquiridos pelo uso de redes sociais e emails.
Aspectos textuais : Quanto a esses aspectos a correção também classificou os textos em cinco níveis ( zero a quatro ), em função da quantidade de erros apresentados.
O desempenho dos estudantes ficou concentrado, majoritariamente, nos níveis 3 e 2, devido à grande ocorrência de problemas de estruturação textual.
Foi muito baixo o número de textos que não apresentaram qualquer problema estrutural e, portanto, ficaram enquadrados no nível 4.
Alguns textos com pior desempenho foram enquadrados no nível 1, enquanto o nível zero foi atribuído a pouquíssimos casos, já que ele revela a existência de um texto sem articulação e com comprometimento do sentido.
Observou -se que a grande maioria dos estudantes não distribuiu as ideias em parágrafos, talvez devido ao pequeno número de linhas disponibilizadas para a resposta da questão ou, quem sabe, pela suposição de que não seria necessária essa divisão por não se tratar de um texto no modelo de uma redação dissertativo-argumentativa, como solicitado nos vestibulares.
Outro aspecto observado na estruturação textual foi a divisão em dois itens, provavelmente motivados pelo encaminhamento do enunciado das questões, que apresentavam dois pontos a serem detalhados : a questão 1 solicitava que os alunos dissertassem sobre as consequências do transporte motorizado ( a ) e apresentassem ações de intervenção por parte do poder público ( b ) ; a questão 2 solicitava que os alunos dissertassem sobre as causas da violência urbana ( a ) e os fatores para evitá -la ( b ).
Essa competência pode ser considerada como a mais problemática no que diz respeito ao desempenho linguístico dos estudantes, porque são muitos os problemas observados, desvios acumulados durante toda a formação escolar e que não se resolvem com um estudo autodidata, como acontece com regras ortográficas ou morfossintáticas : sequência justaposta de ideias sem encaixamentos sintáticos ; redução drástica de estruturas subordinadas, ao lado do aumento na frequência de estruturas coordenadas e absolutas ; redução no uso de conectores para expressar relações lógicas essenciais à construção do texto, substituídas pela exigência de inferência por parte do interlocutor para suprir a sua ausência ; emprego equivocado de operadores que não estabelecem relações lógicas coerentes entre ideias do texto ; emprego inadequado do pronome relativo ( com omissão da preposição ou a utilização de pronome inadequado, como “ onde ” ) ; repetição exaustiva de termos sem a utilização de procedimentos mais sofisticados de substituição ( hiperonímias, hiponímias, nominalizações, expressões metafóricas ) ; frases fragmentadas que comprometem a estrutura lógico-gramatical ; frases formadas apenas por oração subordinada, sem oração principal.
Um importante aspecto a destacar é o baixíssimo desempenho de uma parte dos estudantes em relação à estrutura formal do texto produzido, o que é extremamente preocupante ao se levar em conta que são graduandos em fase final de formação.
São frequentes os casos de desvios de estruturação frasal, com uso inadequado ou ausência de conectivos entre parágrafos e entre frases.
Em uma parte dos textos, falta um mínimo de textualidade e de domínio do registro padrão da língua.
Na verdade, observam -se relações linguísticas quase agramaticais, como as estabelecidas pela sequência de gerúndios sem o apoio de um ponto de partida para a organização das informações gramaticais e semânticas, ou seja, sem uma oração principal.
Quanto à utilização dos mecanismos de referenciação, deve -se destacar a ocorrência, em uma boa parte dos textos, de repetições de palavras ou expressões sem a utilização de termos sinônimos ou pronomes, como seria adequado.
Quanto à utilização dos sinais de pontuação, observou -se uma grande precariedade nos textos analisados.
É muito frequente a ocorrência de parágrafos sem marca interna de pontuação para separar os períodos.
Vale observar que não foi penalizada a ausência de vírgula para destacar locuções ou adjuntos adverbiais de pequena extensão deslocados de posição na frase, por ser um uso opcional.
São os seguintes os tipos de problemas encontrados :  vírgula : utilização de vírgula para separar o sujeito e o predicado ; ocorrência de apenas uma das vírgulas para separar uma palavra, uma expressão ou uma oração encaixada ; uso de vírgula no lugar do ponto para separar ideias que constituem períodos distintos ; ausência de vírgula para separar elementos de uma enumeração ; ausência de vírgula para separar oração adjetiva explicativa ou utilização inadequada para separar oração adjetiva restritiva ;  ponto e vírgula : utilização do ponto e vírgula no lugar de vírgula ;  ponto final : ausência de ponto final para separar períodos.
Aspectos morfossintáticos e vocabulares : Da mesma forma que nos aspectos anteriores os textos foram classificados em níveis de zero a quatro.
O desempenho dos estudantes ficou concentrado, majoritariamente, nos níveis 3 e 2, devido à grande ocorrência de problemas de regência e concordância.
O nível 4 foi atribuído a um número menor de textos.
Alguns textos com pior desempenho foram enquadrados no nível 1, enquanto o nível zero foi atribuído a pouquíssimos casos, já que ele revela a existência de um texto sem o respeito às mínimas exigências morfossintáticas da norma-padrão e com comprometimento do sentido.
Os resultados são muito transparentes em relação aos aspectos mais problemáticos do desempenho dos estudantes.
O desvio mais frequente, em relação à regência, é a falta do sinal indicativo da crase – isso revela que o usuário não tem consciência de que, sob a forma do termo “ a ”, existe a presença de uma preposição “ a ”, exigida pela regência do termo anterior.
Embora em outros exames, como o Enem, a falta de crase seja penalizada em acentuação, nesta correção esse desvio foi considerado no âmbito dos aspectos morfossintáticos.
Outro problema relacionado à regência verbal e à nominal, encontrado frequentemente nas questões, foi a ausência de preposição antes de pronome relativo, processo generalizado na modalidade oral da língua, em situações de registro informal.
Apesar da possibilidade de que essa alteração de regência se generalize no padrão escrito da Língua Portuguesa, como já \textbf{está} ocorrendo até em textos jornalísticos, o não emprego da preposição foi penalizado neste processo de avaliação.
Outro desvio muito frequente diz respeito aos processos de concordância verbal e de concordância nominal.
Quanto à concordância de número, observou -se ausência de marca ( com sujeito anteposto ou posposto ) ou uso indevido ( uso inadequado da marca de plural comandado pelo núcleo plural da locução adjetiva, apesar de o substantivo que funciona como núcleo do sintagma nominal \textbf{estar} no singular ).
Uma ocorrência que se destacou foi a ausência de acento circunflexo na forma plural do presente do indicativo dos verbos “ ter ” e “ vir ”, que foi considerada como um desvio na concordância verbal e não na acentuação gráfica.
Quanto à concordância de gênero, vários casos foram observados, normalmente no âmbito de sintagmas nominais longos, em que o adjetivo \textbf{está} afastado do substantivo.
Deve -se destacar uma ocorrência não observada no ENADE/2013 : o aparecimento da marca de plural em verbos ou adjetivos relacionados a núcleos substantivos no singular, evidenciando um processo de hipercorreção.
Quanto à questão da colocação pronominal, foram poucos os casos observados.
Apesar de serem aspectos relacionados à oralidade, concluiu -se que, no registro escrito formal, a maioria dos estudantes já incorporou regras como a não introdução da frase por um pronome oblíquo e a próclise na presença de um termo atrator.
Não se adotou, entretanto, o padrão excessivamente formal descrito pelas gramáticas normativas em relação à posição do pronome oblíquo em locuções verbais, já que esse uso \textbf{está} muito distante da prática cotidiana, até em textos mais formais.
Quanto aos aspectos vocabulares, alguns tipos de inadequação foram observados : expressões da oralidade apareceram em algumas respostas, mas sem maior relevância do ponto de vista quantitativo ; seleção vocabular incompatível com o contexto, gerando situações de falta de inteligibilidade ; falta de domínio de vocabulário mais abstrato e de maior complexidade, essencial ao desenvolvimento do texto de base dissertativa.
Várias marcas de oralidade foram identificadas, embora não em alta frequência : o uso do pronome relativo “ onde ” como relativo universal, falta de artigo definido antes de substantivo, repetição de palavras por falta de vocabulário, reduções como “ \textbf{tá} ”, “ pra ”, “ pro ”, “ prum ”, expressões informais. concordância de gênero, vários casos foram observados, normalmente no âmbito de sintagmas nominais longos, em que o adjetivo \textbf{está} afastado do substantivo.
Deve -se destacar uma ocorrência não observada no ENADE/2013 : o aparecimento da marca de plural em verbos ou adjetivos relacionados a núcleos substantivos no singular, evidenciando um processo de hipercorreção.
Quanto à questão da colocação pronominal, foram poucos os casos observados.
Apesar de serem aspectos relacionados à oralidade, concluiu -se que, no registro escrito formal, a maioria dos estudantes já incorporou regras como a não introdução da frase por um pronome oblíquo e a próclise na presença de um termo atrator.
Não se adotou, entretanto, o padrão excessivamente formal descrito pelas gramáticas normativas em relação à posição do pronome oblíquo em locuções verbais, já que esse uso \textbf{está} muito distante da prática cotidiana, até em textos mais formais.
Quanto aos aspectos vocabulares, alguns tipos de inadequação foram observados : expressões da oralidade apareceram em algumas respostas, mas sem maior relevância do ponto de vista quantitativo ; seleção vocabular incompatível com o contexto, gerando situações de falta de inteligibilidade ; falta de domínio de vocabulário mais abstrato e de maior complexidade, essencial ao desenvolvimento do texto de base dissertativa.
Várias marcas de oralidade foram identificadas, embora não em alta frequência : o uso do pronome relativo “ onde ” como relativo universal, falta de artigo definido antes de substantivo, repetição de palavras por falta de vocabulário, reduções como “ \textbf{tá} ”, “ pra ”, “ pro ”, “ prum ”, expressões informais. 3.3.2 Componente de Conhecimento Específico Na parte da prova relativa às questões discursivas do Componente de Conhecimento Específico ( Tabela 3.15 ), observa -se que a média foi bem mais baixa do que para as questões discursivas do Componente de Formação Geral.
Enquanto no Componente de Formação Geral a média para estudantes de Tecnologia em Automação Industrial de todo o Brasil foi 49,2, na parte de Conhecimento Específico a média foi 33,0.
A maior média deste componente foi obtida pelos estudantes da região Sul ( 35,7 ), e a menor, pelos da região Centro-Oeste ( 28,2 ).
Quanto à variabilidade das notas, o desvio padrão de todo o Brasil foi 18,7.
O maior desvio padrão foi encontrado na região Sudeste ( 18,8 ), e o menor, na região Norte ( 16,9 ).
A maior nota máxima foi obtida na região Sudeste ( 90,0 ), enquanto a menor nota máxima foi encontrada na região Norte ( 68,3 ).
Além disso, a nota mínima ( 0,0 ) foi obtida por pelo menos um aluno em todas as regiões do Brasil, exceto na região Norte ( 13,3 ).
A mediana do Brasil como um todo foi 33,3.
Nas regiões, a mediana foi : 35,0 na região Norte ; 32,5 na região Nordeste ; 31,7 na região Sudeste ; 36,7 na região Sul ; e 26,7 na região Centro-Oeste.
O Gráfico 3.10 representa a distribuição das notas nas questões discursivas no Componente de Conhecimento Específico.
A moda desta distribuição ocorre no intervalo ( 20 ; 30 ].
Nota -se que a distribuição é decrescente a partir do intervalo modal.
A análise de \textbf{cada} uma destas questões será feita a seguir. 3.3.2.1 Análise da Questão Discursiva 3 do Componente de Conhecimento Específico Na questão 3, cujos resultados aferidos encontram -se descritos na Tabela 3.16.
O desempenho dos estudantes de todo o Brasil nesta questão foi o mais alto dentre as três questões discursivas desse componente, e a média dos estudantes de todo o Brasil foi 56,5.
A menor média nessa questão foi obtida pelos alunos da região Centro-Oeste ( 51,1 ), enquanto a maior média foi obtida nas regiões Norte e Sul ( 57,3 ).
Quanto à variabilidade das notas, o desvio padrão de todo o Brasil foi 28,1.
O maior desvio padrão foi obtido na região Centro-Oeste ( 29,8 ), enquanto o menor foi obtido na região Norte ( 19,0 ).
A nota máxima, 100,0 pontos, foi alcançada por, pelo menos, um aluno de quase todas as regiões, com exceção da região Norte, onde a nota máxima foi 90,0.
A mediana do Brasil como um todo foi 60,0, a mesma nas regiões Nordeste e Sudeste.
A maior mediana nesta questão foi na região Sul ( 70,0 ), e nas regiões Norte e Centro- Oeste a mediana foi 50,0.
A nota mínima ( 0,0 ) foi obtida por, pelo menos, um estudante em todas as regiões do Brasil, exceto na região Norte ( 30,0 ).
O Gráfico 3.11 mostra a distribuição das notas na questão discursiva 3, do Componente de Conhecimento Específico da Área de Tecnologia em Automação Industrial.
Essa distribuição é unimodal com moda nas notas situadas no intervalo ( 60 ; 70 ]. 3.3.2.2 Comentários sobre as respostas à Questão Discursiva 3 A questão avaliava conhecimentos pertinentes aos conteúdos, competências e habilidades presentes nas Diretrizes Curriculares de Área, exigindo nível de profundidade, conforme previsto pelo padrão de respostas compatível para um formando.
Moldando -se à estrutura do enunciado, o padrão de respostas apresentou modelos de solução para os dois pontos focalizados pelo comando da questão : o primeiro ( item ‘ a ’ ), que se detinha sobre os equipamentos de proteção individual a serem utilizados por um eletricista na situação apontada ; e o segundo ( item ‘ b ’ ), que se voltava aos procedimentos de proteção coletiva a serem adotados pelo eletricista em tal situação.
O enunciado era claro e continha informações suficientes para solução da questão.
Quanto ao grau de dificuldade, essa questão pode ser considerada fácil.
Assim, pode -se considerar que a questão poderia ser respondida em pouco tempo e que o espaço reservado para o desenvolvimento da resposta era adequado.
Quanto à relação entre o padrão de respostas e as respostas dadas pelos estudantes, para \textbf{cada} um dos dois itens, destacam -se os aspectos mais relevantes a seguir.
No item ‘ a ’ - alguns estudantes mencionaram EPIs diferentes das que constavam no padrão de respostas.
Como o padrão de respostas esperado priorizou quatro EPIs, só se atribuíram pontos aos mencionados dentro daquele conjunto.
Para o item ‘ b ’ - O padrão de respostas enumerou as quatro ações e só foi atribuída pontuação máxima a quem mencionou as quatro.
Destaque -se que poucos foram os estudantes que acertaram as quatro ações.
A média das notas atribuídas a essa questão ficou no intervalo ( 50 : 60 ].
Em um pouco mais de 10 \% das provas, a questão foi entregue em branco.
Dos estudantes que responderam à questão, cerca de 2 \% receberam nota zero.
A maioria recebeu notas classificadas como medianas, e a um bom número de estudantes foi atribuído notas consideradas boas.
Em geral, no item ‘ a ’, houve uma porcentagem de acertos bastante alta.
A única discrepância com relação ao padrão de respostas ocorreu porque alguns EPIs citados não constavam da lista priorizada pela Comissão.
No item ‘ b ’, poucos foram os estudantes que acertaram as quatro ações listadas.
A ação de desligamento da chave de energia foi bastante citada, mas o travamento da chave de energia teve baixa incidência de citações.
Pelas respostas encontradas, constata -se que há estudantes conhecedores do uso dos EPIs, embora o conhecimento demonstrado quanto aos procedimentos a serem adotados por eletricistas tenha sido apenas regular.
Das questões discursivas do componente específico esta foi certamente a mais fácil. 3.3.2.3 Análise da Questão Discursiva 4 do Componente de Conhecimento Específico A Tabela 3.17 contém as informações relativas à questão 4 do conjunto de questões do Componente de Conhecimento Específico.
O desempenho dos estudantes nessa questão foi inferior ao das questões 3 e 5.
A média geral do Brasil foi 12,3, sendo a menor média registrada na região Centro-Oeste ( 6,1 ), e a maior, na região Norte ( 18,2 ).
A nota máxima ( 100,0 ) foi atingida por, pelo menos, um concluinte em quase todas as regiões do Brasil, exceto nas regiões Norte ( 65,0 ) e Nordeste ( 80,0 ).
A mediana em todo o Brasil foi 0,0, o mesmo valor foi obtido em todas as regiões do Brasil, exceto na região Nordeste ( 5,0 ).
Isso indica que, excetuando -se a região Nordeste, pelo menos 50 \% do total de alunos tirou nota zero nessa questão.
As notas mínimas foram zero em todas as regiões, sem exceção.
O Gráfico 3.12 representa a distribuição de notas da questão discursiva 4, do Componente de Conhecimento Específico.
Essa distribuição é unimodal com moda na categoria “ em branco ”.
Observa -se que uma proporção maior de estudantes deixou essa questão em branco, comparando -se com o ocorrido nas demais questões discursivas de Conhecimento Específico.
O percentual de notas zero também é maior do que nas questões 3 e 5. 3.3.2.4 Comentários sobre as respostas à Questão Discursiva 4 Os conhecimentos exigidos e o nível de profundidade apresentados no padrão de respostas da questão eram pertinentes e compatíveis com os conteúdos, as competências e as habilidades apontados nas Diretrizes Curriculares da Área.
O enunciado \textbf{estava} claro e havia suficiência de informações para a solução da questão.
Essa questão foi considerada difícil, mas para estudantes com conhecimento do assunto, poderia ser resolvida em pouco tempo.
Desse modo, o tempo destinado ao desenvolvimento da resposta à questão foi adequado, assim como também o foi o espaço oferecido para tal.
A questão \textbf{estava} dividida em três itens.
Muitos acertaram o item ‘ a ’, em que se pedia aos estudantes que identificassem as linhas irrelevantes.
Um número razoável Entendeu como linhas irrelevantes aquelas em que o botão que habilita o sistema \textbf{está} desligado.
No item ‘ b ’, um número pequeno de estudantes conseguiu escrever a expressão algébrica mais simplificada do controle de \textbf{cada} eletroválvula.
Por fim, observou -se que muitos dos estudantes que acertaram o item ‘ b ’ acertaram também o item ‘ c ’.
O resultado da correção levou a concluir que a questão deve considerada muito difícil para os estudantes formando na Área.
A média das notas foi baixa.
Registrou -se um elevado número de respostas em branco, aproximadamente 35 \%, e de notas entre zero e 30, consideradas fracas, superou 50 \%.
Menos frequentes foram as respostas classificadas como medianas, com notas entre 35 e 70, e ainda menos frequentes, as consideradas boas, com notas maiores do que 70.
Como já mencionado anteriormente, o item ‘ a ’ foi o que teve uma porcentagem de acertos mais alta.
Os itens ’ b ’ e ‘ c ’ tiveram um índice de acerto baixo, mas muitos estudantes que acertaram o item ‘ b ’ acertaram também o item ‘ c ’.
Embora tratando de temas relevantes para o exercício da profissão, essa foi a questão em que os estudantes mais apresentaram dificuldade.
Os resultados obtidos revelam claramente o pouco aprendizado dos estudantes no assunto relacionado à tabela verdade e linguagem ladder. 3.3.2.5 Análise da Questão Discursiva 5 do Componente de Conhecimento Específico A Tabela 3.18 contém as informações relativas à questão 5 do conjunto do Componente de Conhecimento Específico.
A nota média dos estudantes de todo o Brasil foi 30,3.
A maior média foi registrada na região Sul ( 35,8 ), enquanto a menor média foi registrada na região Centro-Oeste ( 27,5 ).
Quanto à variabilidade das notas, o desvio padrão dos alunos do Brasil, como um todo, foi 28,5.
Enquanto o maior desvio foi encontrado na região Norte ( 30,9 ), o menor foi encontrado na região Sul ( 27,9 ).
A nota mínima ( 0,0 ) é a mesma em todas as regiões do Brasil, sem exceção.
A mediana do Brasil, como um todo, foi 30,0.
A maior mediana foi obtida na região Sul ( 40,0 ), sendo 25,0 nas regiões Norte e Sudeste, e 20,0 nas regiões Nordeste e Centro-Oeste.
Para o conjunto de alunos de Tecnologia em Automação Industrial do Brasil, a nota máxima foi 100,0, sendo obtida por, pelo menos, um estudante em quase todas as regiões do Brasil, exceto nas regiões Norte e Centro-Oeste ( 90,0 em ambas ).
O Gráfico 3.13 apresenta a distribuição das notas da questão discursiva 5 do Componente de Conhecimento Específico.
Destaca -se, o grande número de estudantes que deixaram esta questão em branco, correspondendo à moda da distribuição. 3.3.2.6 Comentários sobre as respostas à Questão Discursiva 5 Os conhecimentos exigidos na questão são pertinentes aos conteúdos, competências e habilidades presentes nas Diretrizes Curriculares da Área.
Da mesma forma, o nível de profundidade apresentado no padrão de respostas era adequado à questão e ao que se espera de um formando.
Comparativamente às outras questões discursivas do componente específico, esta questão pode ser considerada de média dificuldade.
O enunciado \textbf{estava} claro e havia suficiência de informações para solução da questão, que poderia ser elaborada em pouco tempo.
O espaço fornecido para o desenvolvimento da resposta era suficiente.
Houve grande incidência de respostas incompletas, e alguns estudantes mencionaram no item ‘ a ’ ações que deveriam constar no item ‘ b ’.
Verificaram -se que cerca de um quarto das respostas foram entregues em branco.
Dos estudantes que apresentaram respostas, cerca de 60 \% receberam notas classificadas como fracas, cerca de 30 \% obtiveram notas consideradas medianas e apenas um pouco mais de 10 \% receberam notas consideradas boas, sendo que dessas, apenas três receberam nota máxima.
Muitos estudantes obtiveram sucesso no concernente ao item ‘ a ’, que fazia menção ao acionamento do alarme, e pouquíssimos mencionaram o registro da ocorrência no banco de dados.
A maioria dos estudantes acertou o item ‘ b ’, e poucos acertaram por completo o item ‘ c ’.
Pouquíssimos mencionaram que o operador deveria tomar ciência do problema, sendo que alguns responderam que o operador deveria parar o processo.
Esta foi a questão discursiva do componente específico que, provavelmente, mais tomou tempo dos estudantes que buscaram respondê -la, como foi observado pela extensão das respostas produzidas por eles, apesar de o nível de exigência do padrão esperado de resposta não ser alto.
Os resultados alcançados revelam a existência de uma lacuna no aprendizado dos estudantes no tocante ao conhecimento preciso dos papéis exercidos pelos agentes supervisório, CLP e operador.

% matched lemmas: acidentado, banalização, cada, certo, ciclista, criminalização, desvalorização, estar, lembrar
\end{textsample}
